\chapter{Khoa học là gì?}

\

Với mỗi người, thế giới vật lí đem lại cho họ những ý thức khác nhau. Với các bà bán hàng rong, việc tìm hiểu giá tiền các loại rau củ hay các loại thịt gân, nạc, đùi là thế giới quan của họ. Với các bác nông dân, việc tìm hiểu thời tiết, tìm hiểu các giống cây trồng và phương pháp nuôi gia súc có khi mới là thứ họ quan tâm. Với các người viết nhạc, có thể tâm hồn lại bởi háo hức với những giai điệu, những nốt nhạc, những bản nhạc hay. Tất cả những sự thay đổi của thế giới, từ tự nhiên đến xã hội, đều tác động đến từng con người, Việt Nam hay thế giới. Những người làm khoa học cũng không phải là ngoại lệ. Họ cũng tìm hiểu thế giới xung quanh mình, và cũng hành động trên nó. Vậy, sự khác biệt giữa khoa học và những chuyên ngành khác là gì?

Khi chúng ta không biết nghĩa của một từ, chúng ta thường tra từ điển để hiểu nghĩa của nó. Và bây giờ, chúng ta đang không biết nghĩa của từ ``khoa học'', vậy chúng ta hãy cùng nhau tra từ điển để hiểu nghĩa của nó. Theo \cite{hoangphe2022tudientienviet},
\begin{quotation}
    ``\defText{Khoa học} là hệ thống tri thức tích lũy trong quá trình lịch sử và được thực tiễn chứng minh, phản ánh những quy luật khách quan của thế giới bên ngoài cũng như của hoạt động tỉnh thần của con người, giúp con người có khả năng cải tạo thế giới hiện thực.''.
\end{quotation}

Định nghĩa tương đối dài dòng, do đó, chúng ta sẽ cần phân tích một chút để giải nghĩa từng phần của định nghĩa và từ đó phân biệt khoa học với các lĩnh vực khác.

\section{Tính hệ thống của khoa học}

\

Một mẩu thông tin đúng mà đa số mọi người công nhận với nhau là đúng được gọi là một \defText{khẳng định} hay \defText{sự thật}. Các khẳng định có thể là:
\begin{itemize}
    \item ``Mọi vật có sự sống cuối cùng đều sẽ chết.'';
    \item ``Hình vuông là tứ giác có bốn cạnh bằng nhau và bốn góc vuông.'';
    \item ``Nước đá thì là chất rắn, còn nước ở nhiệt độ thông thường là chất lỏng.''.
\end{itemize}

Trong trường hợp mà mẩu thông tin chưa được công nhận hay được kiểm chứng, thì sao? Nếu chúng ta có thể kiểm tra được tính chính xác của nó, thì chúng ta sẽ gọi nó là một \defText{giả thuyết}. Ví dụ:
\begin{itemize}
    \item ``Có sự sống ngoài Trái Đất.'';
    \item ``Loại thuốc $A$ này chữa được xoang.''.
\end{itemize}
Các phương pháp để kiểm tra các giả thuyết được gọi chung là \defText{phương pháp thực nghiệm khoa học}. Khi tổng hợp đủ nhiều các khẳng định và giả thuyết đã được kiểm chứng, chúng ta sẽ có thể rút ra được những tính chất chung, hay còn gọi là các \defText{quy luật}. Những quy luật mà bạn đọc có thể đã biết đến bao gồm:
\begin{itemize}
    \item \textcolor{colorEmphasisCyan}{Định luật vạn vật hấp dẫn}: Định luật này giải thích quy luật tương tác giữa các vật thể có khối lượng trong vũ trụ. Niu-tơn\footnote{Isaac Newton (1643 -- 1727).} không bỗng nhiên nghĩ ra định luật này chỉ vì một quả táo rơi. Nó là kết quả của sự tổng hợp dữ liệu khổng lồ trước đó. Kép-lơ\footnote{Johannes Kepler (1571 -- 1630).} đã tổng hợp dữ liệu về chuyển động của các hành tinh quan sát bởi Bra-hê\footnote{Tycho Brahe (1546 -- 1601).} để rút ra được ba định luật chuyển động hành tinh, từ đó Niu-tơn mới có thể sử dụng toán học để thống nhất hiện tượng ``vật rơi dưới đất'' và ``hành tinh quay trên trời'' thành một công thức duy nhất. Định luật này đã được kiểm chứng qua vô số thực nghiệm và quan sát, từ việc dự đoán quỹ đạo của các hành tinh đến việc giải thích hiện tượng thủy triều trên Trái Đất.
    \item \textcolor{colorEmphasis}{Định luật bảo toàn và chuyển hóa năng lượng}: Một trong những quy luật cơ bản nhất, khẳng định năng lượng không tự nhiên sinh ra hay mất đi. Tuy quy luật này thường được phát biểu trong phạm trừ của vật lí, nó không xuất phát từ một thí nghiệm duy nhất mà là sự giao thoa của nhiều ngành. Giun\footnote{James Prescott Joule (1818 -- 1889).} đã dùng một quả nặng rơi để làm quay các cánh quạt trong nước và chứng minh rằng công cơ học biến thành nhiệt năng theo một tỉ lệ không đổi. Mai-ơ\footnote{Julius Robert von Mayer (1814 -- 1878).}, sau khi quan sát được sự khác biệt của màu máu giữa các vùng, ông đã đưa ra kết luận rằng thức ăn không chỉ tạo ra nhiệt lượng mà còn tạo ra lực chuyển động và nhiệt giữ ấm cho cơ thể, đóng góp vào việc xây dựng định luật này. 
    
    Định luật bảo toàn và chuyển hóa năng lượng cũng là một trong những nguyên nhân khách quan dẫn đến sự tồn tại của chủ nghĩa duy vật biện chứng của Các Mác\footnote{Karl Marx (1818 -- 1883).} và Ăng-ghen\footnote{Friedrich Engels (1820 -- 1895).}, khi họ sử dụng quy luật này để giải thích sự vận động luân hồi của vật chất.
    \item \textcolor{colorEmphasisGreen}{Định luật Ôm}: Quy luật này mô tả mối quan hệ giữa các thông số cơ bản trong một mạch điện, được kiểm chứng qua vô số thực nghiệm về dòng điện. Không có một kĩ sư điện hay điện tử nào có thể làm việc mà không hiểu về định luật này.
    \item \textcolor{colorEmphasisPurple}{Định luật Men-đen về di truyền học}\footnote{Gregor Mendel (1822 -- 1884).}: Định luật này giải thích cách thức di truyền các đặc điểm từ thế hệ này sang thế hệ khác, dựa trên các thí nghiệm với cây đậu Hà Lan, từ đó tạo nên nền tảng của di truyền học hiện đại.
\end{itemize}